\documentclass[10pt,a4paper]{amsart}
\usepackage[margin=19mm]{geometry}
\usepackage{amsmath,amssymb,mathtools}
\usepackage[scr=rsfs,cal=euler]{mathalfa}
\usepackage{xfrac}
\usepackage[unicode,hidelinks,bookmarks=false]{hyperref}
\newcommand{\ie}{i.e.\ }
\newcommand{\cf}{cf.\ }
\newcommand{\resp}{respectively\ }
\newcommand{\Subsection}{\subsection}
\providecommand{\nobreakdash}{\nobreak\hbox{-}}
\setlength{\emergencystretch}{2em}
\multlinegap0pt
\usepackage{amssymb}
\usepackage{xcolor}
\usepackage{tikz}
\usepackage[shortlabels]{enumitem}
\usepackage{float}
\newtheorem{lemma}{Lemma}
\newtheorem{Proof}{Proof}
\DeclareMathOperator{\N}{\mathbb{N}}
\DeclareMathOperator{\iii}{\mathfrak{i}}
\DeclareMathOperator{\jjj}{\mathfrak{j}}
\DeclareMathOperator{\kkk}{\mathfrak{k}}
\title{Weakly separated self-affine carpets}
\author{Bal\'azs B\'ar\'any, Levente David}
\date{Source: arXiv:2506.06851v2 (CC BY 4.0)}
\begin{document}
\maketitle

\noindent\emph{Source and licence.} Weakly separated self-affine carpets. Version arXiv:2506.06851v2 (CC BY 4.0), distributed by arXiv under the Creative Commons Attribution 4.0 International licence, http://creativecommons.org/licenses/by/4.0/, as recorded in the arXiv licence metadata for this exact version and shown on the abstract page https://arxiv.org/abs/2506.06851v2. Full licence notice: \texttt{licenses/arxiv-2506.06851-CC-BY-4.0.txt}.\par\smallskip

\noindent\textit{Editorial scope.} Counted unit: the article's lemma ex2l2 (source0.tex lines 1623--1629). The article's complete original proof of it is delivered with it (source0.tex lines 1631--1756, one \texttt{proof} environment of 126 lines). No mathematics is written, repaired or re-derived: the statement and the proof are the article's own lines, in the article's own notation, and no retained line is substituted or re-flowed.  Retained uncounted as the source-local context the proof needs: lines 1553--1561.  Nothing was omitted from any retained span, so the package carries no omission record.  No retained line contains a citation, so the wrapper carries no bibliography and no bibliography entry.  No retained line contains a figure environment, an image-inclusion command or drawing code, so no image is delivered and the package declares no asset dependency.  Editorial formatting only, in addition: an AMS \texttt{amsart} preamble that restates the article's own declarations and macros used by the retained lines, namely the environments \texttt{lemma}, \texttt{Proof} on separate counters, together with the standard packages those lines need.  The retained environments print in this document as Lemma 1, Proof 1 in order of appearance, whereas the article prints the same items with the numbering of its own full document.  The complete native source of the article as distributed in the arXiv e-print for this exact version is bundled unchanged as \texttt{sources/179488/source0.tex}.
\begin{lemma} \label{ex2l1}

For any $n \in \N^+$ let 
\begin{align*}
\Xi_n :=&\ \Big\{ \iii \in \Sigma^n \ \Big| \ \forall k \in \{1,\dots,n-1\}: (\iii_k \iii_{k+1}) \notin W \Big\} \\
\Theta_n :=&\ \Big\{ \iii \in \Sigma^n \ \Big| \ \forall k \in \{1,\dots,n-1\}: (\iii_k \iii_{k+1}) \notin V \Big\}.
\end{align*}
Now $\Xi_n$ satisfies the definition of $\Gamma_2^{\{n\}}$ and $\Theta_n$ satisfies the definition of $\Sigma^{\{n\}}$.
\end{lemma}

\begin{lemma} \label{ex2l2}
For $\iii \in \Xi_n$ and $j \in \Sigma$ such that $j\iii_1 \notin W$,
\begin{align*}
R_{(j \iii)} = \begin{cases}(\frac{k+2}{k+1})^\alpha \cdot R_{\iii}& \text{if } j = 3 \text{ and } \iii = 3^k4\kkk \text{ for } k \leq n-1,\, \kkk \in \Xi_{n-k-1}, \\ 
R_{\iii} & \text{else}. \end{cases}
\end{align*}
\end{lemma}

\begin{Proof}

By Lemma \ref{ex2l1}, for any $\iii \in \Xi_n$
\[ R_{\iii} = (\#\{S\in\mathbb{G}_n\, \big|\, \mathrm{p}_{2} S = \mathrm{p}_2 S_{\iii}\})^\alpha = (\#\{\jjj \in \Theta_n\, \big|\, \mathrm{p}_{2}(\jjj) = \iii\})^\alpha. \]

We make some observations first. 
Let us fix $j \in \{1,2\}$. Then we have the following two facts:
\begin{itemize}
\item For any $\ell \in \Sigma$ such that $j\ell \in \Theta_2$, $\{(ab)\in \Theta_2| \mathrm{p}_2(ab) = \mathrm{p}_2(j\ell) \} = \{j\ell\}$. Thus, for any $\iii \in \Xi_{n-1}$ such that $j\iii \in \Xi_n$
\begin{align} \label{eqex2ob1}
\big\{\jjj \in \Theta_{n}\, \big|\, \mathrm{p}_2(\jjj) = \mathrm{p}_2(j\iii)\big\} = \big\{j\jjj \in \Theta_{n}\, \big|\, \mathrm{p}_2(\jjj) = \mathrm{p}_2(\iii)\big\}.
\end{align}
In particular, for every $\iii\in\Xi_{n-1}$ such that $j\iii\in\Xi_n$
\[R_{j\iii} = R_{\iii}.\]
\item Let $j \neq k\in\Sigma$. If $\mathrm{p}_2(kc) = \mathrm{p}_2(ab)$ where $(kc), (ab) \in \Theta_2$, then $a \neq j$. Thus, if $k\iii, \jjj \in \Theta_n$ are such that $\mathrm{p}_2(k\iii) = \mathrm{p}_2(\jjj)$, then
\begin{align}\label{eqex2ob2}
\jjj_1 \neq j.
\end{align}
\end{itemize}

Let $\iii \in \Xi_{n-1}$ such that $4\iii \in \Xi_n$. Now, $\iii_1 \neq 1$ as $41 \in W$. Thus, by \eqref{eqex2ob2}
\[
\big\{\jjj \in \Theta_{n}\, \big|\, \mathrm{p}_2(\jjj) = \mathrm{p}_2(4\iii)\big\} = \big\{4\jjj \in \Theta_{n}\, \big|\, \mathrm{p}_2(\jjj) = \mathrm{p}_2(\iii)\big\}.
\]
Therefore,
\[R_{4\iii} = R_{\iii}.\]

Let $k\in \N, n\geq k$ and $\iii \in \Theta_{n-k}$ such that $\iii_1 = 2$. Then $3^k\iii \in \Theta_n$. Then, $3^k\iii = 3^k2\dots$ and since $\{\jjj \in \Theta_{k+1} | \mathrm{p}_2(\jjj) = \mathrm{p}_2(3^k2)\} = \{3^k2\}$ by \eqref{eqex2ob2}
\[R_{3^k\iii} = R_{\iii}.\]

Finally, let $\iii = 3^k4\kkk \in \Xi_n$, where $k \leq n-1$. Then one easily has $\{\jjj \in \Theta_{k+1}| \mathrm{p}_2(\jjj) = \mathrm{p}_2(3^k4)\} = \big\{3^\ell41^{k-\ell} \, \big|\, \ell \in \{0,1,\dots,k\}\big\}$ and since $3^k4\kkk \in \Xi_n$, either $\kkk$ is the empty word, or $\kkk_1 \neq 1$. Thus, by \eqref{eqex2ob2}
\[
\#\{\jjj \in \Theta_n\, \big|\, \mathrm{p}_{2}(\jjj) = \mathrm{p}_2 (3^k4\kkk)\} = \#\{\jjj \in \Theta_{k+1}\, \big|\, \mathrm{p}_{2}(\jjj) = \mathrm{p}_2 (3^k4)\} \cdot \#\{\jjj \in \Theta_{n-k-1}\, \big|\, \mathrm{p}_{2}(\jjj) = \mathrm{p}_2 (\kkk)\}.
\] 
Therefore, for every $k\geq 1$
\[R_{3^k4\iii} = (k+1)^\alpha \cdot R_{\iii}=\frac{(k+1)^\alpha}{k^\alpha} k^\alpha R_{\iii}=\frac{(k+1)^\alpha}{k^\alpha} R_{3^{k-1}4\iii}.\]


%
%then $R_{\iii} = (k+1)^\alpha R_{\kkk}$. Indeed, 
%\[ (R_{3\iii})^{1/\alpha} = \#\{S\in\mathbb{G}_n\, \big|\, \mathrm{p}_{2} S = \mathrm{p}_2 S_{\iii}\} = \#\{\jjj \in \Theta_n\, \big|\, \mathrm{p}_{2}(\jjj) = \mathrm{p}_2 (3^k4\kkk)\} = (k+1) \cdot \#\{\jjj \in \Theta_n\, \big|\, \mathrm{p}_{2}(\jjj) = \mathrm{p}_2 (\kkk)\}, \]
%where for the last equality, we used that 
%\[ \big\{\jjj \in \Xi_{k+1}\, \big|\, \mathrm{p}_2(\jjj) = \mathrm{p}_2(3^k4)\big\} = \big\{3^\ell41^{k-\ell} \, \big|\, \ell \in \{0,1,\dots,k\}\big\}\]
%and that
%\[
%\big\{\jjj \in \Xi_{n-k}\, \big|\, \mathrm{p}_2(\jjj) = \mathrm{p}_2(4\kkk)\big\} = \big\{4\jjj \in \Xi_{n-k}\, \big|\, \mathrm{p}_2(\jjj) = \mathrm{p}_2(\kkk)\big\}.
%\]
%
%Then we show 
%
%
%\color{blue}
%Let $\iii \in \Xi_n$ and $j \in \{1,2,4\}$ be such that $j\iii \in \Xi_n$. By Lemma \ref{ex2l1},
%\begin{align*}
%\#\{S\in\mathbb{G}_{n+1}\, \big|\, \mathrm{p}_{2} S &= \mathrm{p}_2 S_{j\iii}\} = \#\{\jjj \in \Theta_{n+1}\, \big|\, \mathrm{p}_{2}(\jjj) = \mathrm{p}_2 (j\iii)\}.
%\end{align*}
%Now, simce $j\iii \in \Xi_n$, we have $j \iii_1 \notin W$. This means that we are counting functions whose projection does not have an equivalent banned name, and thus has a unique row. Thus
%\begin{align*}
%\#\{\jjj \in \Theta_{n+1}\, \big|\, \mathrm{p}_{2}(\jjj) = \mathrm{p}_2 (j\iii)\} = \#\{\jjj \in \Theta_{n}\, \big|\, \mathrm{p}_{2}(\jjj) = \mathrm{p}_2 (\iii)\}.
%\end{align*}
%Whence, 
%\[R_{j\iii} = R_{\iii}.\]
%
%Let $\iii \in \Xi_n$ be such that $s(\iii) \in \{2,3\}$. By Lemma \ref{ex2l1},
%\begin{align*}
%\#\{S\in\mathbb{G}_{n+1}\, \big|\, \mathrm{p}_{2} S &= \mathrm{p}_2 S_{3\iii}\} = \#\{\jjj \in \Theta_{n+1}\, \big|\, \mathrm{p}_{2}(\jjj) = \mathrm{p}_2 (3\iii)\}.
%\end{align*}
%If $s(\iii) = 2$, then 
%\[
%\#\{\jjj \in \Theta_{n+1}\, \big|\, \mathrm{p}_{2}(\jjj) = \mathrm{p}_2 (3\iii)\} = \#\{\jjj \in \Theta_{n+1}\, \big|\, \mathrm{p}_{2}(\jjj) = \mathrm{p}_2 (32\kkk)\} = \#\{\jjj \in \Theta_{n+1}\, \big|\, \mathrm{p}_{2}(\jjj) = \mathrm{p}_2 (32\kkk)\},
%\]
%since $32 \notin W$. If $s(\iii) = 3$, then $\iii = 3\dots3\kkk$, where $\kkk$ is either in $\Xi_*$ and starts with $2$, or is the empty word. In both cases, 
%\[
%\#\{\jjj \in \Theta_{n+1}\, \big|\, \mathrm{p}_{2}(\jjj) = \mathrm{p}_2 (33\dots3\kkk)\} = \#\{\jjj \in \Theta_{n+1}\, \big|\, \mathrm{p}_{2}(\jjj) = \mathrm{p}_2 (33\dots3\kkk)\} = \#\{\jjj \in \Theta_{n+1}\, \big|\, \mathrm{p}_{2}(\jjj) = \mathrm{p}_2 (33\dots3\kkk)\},
%\]
%
%Let $\iii \in \Xi_n$ be such that $s(\iii) = 4$. Thus $\iii = 4\kkk$ for $\kkk \in \Xi_{n-1}$. Notice that $34 \notin W$, thus $3\iii \in \Xi_n$. By Lemma \ref{ex2l1},
%\begin{align*}
%\#\{S\in\mathbb{G}_{n+1}\, \big|\, \mathrm{p}_{2} S &= \mathrm{p}_2 S_{3\iii}\} = \#\{\jjj \in \Theta_{n+1}\, \big|\, \mathrm{p}_{2}(\jjj) = \mathrm{p}_2 (34\kkk)\}.
%\end{align*}
%Since $\{ \iii \in \Theta_2\,|\, \mathrm{p}_2(\iii) = \mathrm{p}_2(34)\} = \{34,41\}$,
%\begin{multline*}
%\#\{\jjj \in \Theta_{n+1}\, \big|\, \mathrm{p}_{2}(\jjj) = \mathrm{p}_2 (34\kkk)\} = \#\{\jjj \in \Theta_{n+1}\, \big|\, \jjj_1 = 3,\, \mathrm{p}_{2}(\jjj|_{(1,n]}) = \mathrm{p}_2 (4\kkk)\} \\ + \# \{\jjj \in \Theta_{n+1}\, \big|\, \jjj_1 = 4,\, \mathrm{p}_{2}(\jjj|_{(1,n]}) = \mathrm{p}_2 (1\kkk)\}.
%\end{multline*}
%If $\mathrm{p}_2(ij) = \mathrm{p}_2(4\kkk_1)$ then $i \in \{3,4\}$. Thus $3i \notin \{21,31\}$. Therefore
%\begin{align*}
%\#\{\jjj \in \Theta_{n+1}\, \big|\, \jjj_1 = 3,\, \mathrm{p}_{2}(\jjj|_{(1,n]}) = \mathrm{p}_2 (4\kkk)\} = \#\{\jjj \in \Theta_{n}\, \big|\, \mathrm{p}_{2}(\jjj) = \mathrm{p}_2 (4\kkk)\}.
%\end{align*}
%Similarly,
%\begin{align*}
%\# \{\jjj \in \Theta_{n+1}\, \big|\, \jjj_1 = 4,\, \mathrm{p}_{2}(\jjj|_{(1,n]}) = \mathrm{p}_2 (1\kkk)\} = \# \{\jjj \in \Theta_n\, \big|\, \mathrm{p}_{2}(\jjj) = \mathrm{p}_2 (1\kkk)\}.
%\end{align*}
%Since $i \notin \{1,4\}$ for any $(ij) \in V$,
%\begin{align*}
%\#\{\jjj \in \Theta_{n}\, \big|\, \mathrm{p}_{2}(\jjj) = \mathrm{p}_2 (4\kkk)\} = \# \Theta_{n-1} = \#\{\jjj \in \Theta_{n}\, \big|\, \mathrm{p}_{2}(\jjj) = \mathrm{p}_2 (1\kkk)\}.
%\end{align*}
%Therefore, we can conclude that
%\begin{align*}
%R_{3\iii} = 2^\alpha \cdot \#\{\jjj \in \Theta_{n}\, \big|\, \mathrm{p}_{2}(\jjj) = \mathrm{p}_2 (4\kkk)\} = 2^\alpha \cdot R_{\iii}.
%\end{align*}
%
%Finally, let $\iii \in \Xi_n$ be such that $s(\iii) = \underbrace{33\dots33}_k4$ for $k \geq 1$. Then 
%
%Let $\iii \in \Xi_n, \jjj, \kkk \in \Psi$. A similar proof to the proof of Lemma \ref{ex2l1} shows that 
%\begin{align*}
%\Upsilon_n := \Big\{ \iii \in \{1,2,3,4\}^n \ \Big| \ \forall k \in \{1,\dots,n-1\}: (\iii_k \iii_{k+1}) \notin \{21,31\} \Big\}
%\end{align*}
%satisfies the definition of $\Sigma^{\{n\}}$. Thus, we only need to work with
%\[ \mathrm{p}_2 (34) = \mathrm{p}_2 (41) \text{ while } S_{34} \neq S_{41}. \]
%
%From Lemma \ref{ex2l1} it is clear that we can look at the evolution of the row structure (in particular, ), which can be described with pairs of the elements from $\Sigma$. 
%
%
%Let $V:= \{14,24,34\}$ be the allowed pairs that generate the functions for the elements of $W$. We have that
%\begin{align*}
%R_{(\iii \jjj \kkk)} = R_{(\iii \jjj)}
%\end{align*}
%if $(\jjj_{|\jjj|}\kkk_1) \neq V$.

%For $\iii \in \Sigma^*$, let $\iii = (\jjj_1, \jjj_2, \dots, \jjj_{n(\iii)})$ be its minimal (w.r.t $n$) description with $\jjj_k \in \Phi$. We call this the $\Phi$-description of $\iii$. Let us denote by $\Phi_n := \{\jjj \in \Sigma^* \,|\, n(\iii) = n\}$. Now, we prove by induction on $n$.
%
%For $n=1$, $\iii \in \Phi = \Phi_1$ gives $R_{\iii} = 1$.
%
%If the recurrence relation holds for $n$, then let $\iii \in \Sigma^*$, let $\iii = (\jjj_1, \jjj_2, \dots, \jjj_n)$

\end{Proof}

\end{document}
